% Standalone problem document, generated from the adjacent README.md.
% Compile with XeLaTeX. Mathematical content is maintained in Markdown.
\documentclass[11pt,a4paper]{article}
\usepackage[fontsize=10.5pt]{scrextend}
\usepackage[margin=22mm,headheight=15pt,headsep=9mm,footskip=11mm]{geometry}
\usepackage{fontspec}
\setmainfont{texgyrepagella}[Extension=.otf,UprightFont=*-regular,BoldFont=*-bold,ItalicFont=*-italic,BoldItalicFont=*-bolditalic]
\setsansfont{texgyreheros}[Extension=.otf,UprightFont=*-regular,BoldFont=*-bold,ItalicFont=*-italic,BoldItalicFont=*-bolditalic]
\setmonofont{texgyrecursor}[Extension=.otf,UprightFont=*-regular,BoldFont=*-bold,ItalicFont=*-italic,BoldItalicFont=*-bolditalic,Scale=MatchLowercase]
\usepackage{amsmath,amssymb}
\usepackage{unicode-math}
\setmathfont{texgyrepagella-math.otf}
\usepackage{xcolor,microtype,enumitem,needspace,fancyhdr}
\usepackage{bookmark}
\definecolor{ink}{HTML}{17344B}
\definecolor{accent}{HTML}{197D80}
\definecolor{muted}{HTML}{596773}
\hypersetup{colorlinks=true,linkcolor=accent,urlcolor=accent,pdftitle={MI-05 - The
Marcus--de Oliveira determinantal
conjecture},pdfauthor={Open Problems in Numerical Linear Algebra}}
\urlstyle{same}
\setlength{\parindent}{0pt}
\setlength{\parskip}{5pt plus 1pt minus 1pt}
\linespread{1.04}
\setlength{\emergencystretch}{3em}
\widowpenalty=10000
\clubpenalty=10000
\displaywidowpenalty=10000
\allowdisplaybreaks[1]
\setlist{nosep,leftmargin=1.5em}
\providecommand{\tightlist}{\setlength{\itemsep}{0pt}\setlength{\parskip}{0pt}}
\providecommand{\pandocbounded}[1]{#1}
\pagestyle{fancy}
\fancyhf{}
\fancyhead[L]{\sffamily\footnotesize\color{muted}OPEN PROBLEMS IN NUMERICAL LINEAR ALGEBRA}
\fancyhead[R]{\sffamily\bfseries\footnotesize\color{accent}MI-05}
\fancyfoot[L]{\parbox[b]{0.9\textwidth}{\sffamily\fontsize{8}{10}\selectfont\color{muted}Editorial ratings; a bounded search cannot certify that a problem remains unsolved.\\Literature check: 2026-09-08}}
\fancyfoot[R]{\sffamily\footnotesize\color{muted}\thepage}
\renewcommand{\headrulewidth}{0.3pt}
\renewcommand{\headrule}{\hbox to\headwidth{\color{accent}\leaders\hrule height \headrulewidth\hfill}}
\makeatletter
\renewcommand{\section}{\Needspace{5\baselineskip}\@startsection{section}{1}{0pt}{12pt plus 2pt minus 2pt}{4pt}{\sffamily\bfseries\large\color{ink}}}
\renewcommand{\subsection}{\Needspace{4\baselineskip}\@startsection{subsection}{2}{0pt}{9pt plus 1pt minus 1pt}{3pt}{\sffamily\bfseries\normalsize\color{ink}}}
\makeatother
\setcounter{secnumdepth}{0}
\begin{document}
{\sffamily\small\color{muted}Matrix inequalities and norms\par}
\vspace{3pt}
{\raggedright\hyphenpenalty=10000\exhyphenpenalty=10000\sffamily\bfseries\fontsize{21}{24}\selectfont\color{ink}The
Marcus--de Oliveira determinantal conjecture\par}
\vspace{7pt}
{\sffamily\small\textbf{Difficulty:} extreme\quad\textbf{Importance:} broadly
interesting\par}
\vspace{4pt}
{\color{accent}\hrule height 0.8pt}
\vspace{6pt}
\subsection{Problem statement}\label{problem-statement}

For every integer \(n\ge1\), let \(A,B\in\mathbb C^{n\times n}\) be
normal matrices, meaning \(AA^*=A^*A\) and \(BB^*=B^*B\). List their
eigenvalues with algebraic multiplicity as \(a_1,\ldots,a_n\) and
\(b_1,\ldots,b_n\). Is

\[\det(A+B)\in\operatorname{conv}\left\{\prod_{i=1}^n(a_i+b_{\sigma(i)}):\sigma\in S_n\right\}?\]

Here \(S_n\) is the permutation group and the convex hull is taken in
\(\mathbb C\cong\mathbb R^2\).

\subsection{Why it matters}\label{why-it-matters}

The conjecture would bound a determinant using only two spectra, even
when the matrices are not simultaneously diagonalizable. It is a
fundamental spectral enclosure problem for structured matrix
computations.

\subsection{References}\label{references}

\begin{enumerate}
\def\labelenumi{\arabic{enumi}.}
\tightlist
\item
  A. Kovacec, \emph{A conjecture more precise and stronger than the one
  by Marcus and de Oliveira}, DMUC Preprint 26-11 (7 April 2026), §1,
  Conjectures 1 and 1′.
  \href{https://www.mat.uc.pt/preprints/ps/p2611.pdf}{Primary paper}.
\item
  N. Bebiano and J. P. da Providência, \emph{Revisiting the Marcus--de
  Oliveira conjecture}, Mathematics 13(5) (2025), 711, §1.
  \href{https://doi.org/10.3390/math13050711}{DOI}.
\item
  J. I. Mulero-Martínez, \emph{A variational framework for determinantal
  inequalities of normal matrices: Successes and obstructions}, Linear
  Algebra and its Applications 740 (2026), 19--38; abstract and
  structured-class results.
  \href{https://doi.org/10.1016/j.laa.2026.03.019}{DOI}.
\end{enumerate}

\subsection{Status check --- 2026-09-08}\label{status-check-2026-09-08}

The April 2026 preprint proposes a stronger conjecture rather than
proving this one. The July 2026 journal paper explicitly describes the
unrestricted conjecture as open and proves structured cases. Searches
included \textit{Marcus de Oliveira conjecture 2026 proof},
\textit{Marcus Oliveira determinant counterexample}, and the exact
2026 titles. No complete proof or counterexample was located. Primary
full text was checked for reference 1; reference 3's status evidence was
its publisher abstract, not an independently audited proof.
\end{document}
